% !TEX program = pdflatex
\documentclass[aspectratio=169]{beamer}

\usepackage[T5]{fontenc}
\usepackage[utf8]{inputenc}
\usepackage[english,vietnamese]{babel}
\usepackage{tikz}
\usetikzlibrary{arrows.meta,positioning,shapes.geometric}
\usepackage{booktabs}
\usepackage{tabularx}

\definecolor{lacquer}{HTML}{1A120B}
\definecolor{gold}{HTML}{C69A3F}
\definecolor{cream}{HTML}{EFE6D2}
\definecolor{jade}{HTML}{1D6B5A}

\usetheme{Madrid}
\usecolortheme{default}
\setbeamertemplate{navigation symbols}{}
\setbeamertemplate{footline}[frame number]
\setbeamertemplate{itemize items}[circle]
\setbeamertemplate{enumerate items}[default]

\setbeamercolor{palette primary}{bg=lacquer,fg=cream}
\setbeamercolor{palette secondary}{bg=jade,fg=white}
\setbeamercolor{frametitle}{bg=lacquer,fg=cream}
\setbeamercolor{title}{fg=cream,bg=lacquer}
\setbeamercolor{block title}{bg=gold!85!black,fg=white}
\setbeamercolor{block body}{bg=gold!8,fg=black}
\setbeamercolor{structure}{fg=gold!90!black}

\title[BeVietnam]{BeVietnam: Ứng Dụng Tư Duy Tính Toán Cho Du Lịch Văn Hoá}
\subtitle{Từ pain point của du khách đến feed, bản đồ động, storyline và capture judge}
\author{Khang Phuc Nguyen -- BeVietnam Team}
\institute{Môn học Tư duy tính toán \\ HCMUS -- VNUHCM}
\date{Tháng 06/2026}

\begin{document}

\begin{frame}
  \titlepage
\end{frame}

\begin{frame}{1. Thông điệp chính}
\begin{block}{Điều nhóm muốn chứng minh}
BeVietnam không chỉ là một ứng dụng có AI, mà là một lời giải được xây dựng bằng tư duy tính toán cho một bài toán kinh doanh thật trong du lịch văn hoá.
\end{block}
\begin{itemize}
    \item Bắt đầu từ pain point của du khách.
    \item Phân rã thành các bài toán nhỏ có thể tính toán.
    \item Thiết kế dữ liệu, thuật toán, workflow và fallback.
\end{itemize}
\end{frame}

\begin{frame}{2. Bối cảnh}
\begin{itemize}
    \item Việt Nam có nhiều tài sản du lịch văn hoá: di tích, ẩm thực, lễ hội, kiến trúc, đời sống địa phương.
    \item Nhưng trải nghiệm số hiện nay thường dừng ở bản đồ, ảnh, rating và review.
    \item Du khách biết một nơi nổi tiếng, nhưng chưa chắc hiểu nơi đó có ý nghĩa gì.
    \item Huế được chọn làm pilot vì có mật độ văn hoá cao và phạm vi đủ rõ.
\end{itemize}
\end{frame}

\begin{frame}{3. Vấn đề người dùng}
\begin{enumerate}
    \item Thông tin du lịch bị phân tán.
    \item Thiếu chiều sâu văn hoá trong mô tả địa điểm.
    \item Dịch ngôn ngữ không đồng nghĩa với hiểu văn hoá.
    \item Du khách không biết nên làm gì khi đã đến nơi.
    \item Thời tiết, khoảng cách, giao thông và độ đông thay đổi trải nghiệm.
    \item AI và dịch vụ ngoài có thể lỗi, hệ thống cần fallback.
\end{enumerate}
\end{frame}

\begin{frame}{4. Câu hỏi sản phẩm}
\begin{center}
\Large
\textit{Nên đi đâu bây giờ?\\ Vì sao nơi đó đáng đi?\\ Và khi đến đó nên làm gì?}
\end{center}
\vspace{0.3cm}
\begin{block}{Chuyển thành bài toán tính toán}
Cho người dùng, ngữ cảnh hiện tại và tập địa điểm, hệ thống cần trả về gợi ý, điểm phù hợp, bản đồ ưu tiên và nhiệm vụ văn hoá.
\end{block}
\end{frame}

\begin{frame}{5. Tư duy tính toán được áp dụng}
\begin{itemize}
    \item \textbf{Decomposition}: phân rã bài toán du lịch văn hoá.
    \item \textbf{Pattern Recognition}: nhận diện các mẫu ra quyết định của du khách.
    \item \textbf{Abstraction}: mô hình hoá Place, Context, Fact, Quest, Capture.
    \item \textbf{Algorithm Design}: tính điểm, xếp hạng, chọn nhiệm vụ, xác minh.
    \item \textbf{Evaluation}: đánh giá giá trị người dùng và độ tin cậy hệ thống.
\end{itemize}
\end{frame}

\begin{frame}{6. Decomposition}
\begin{tikzpicture}[
    box/.style={draw=gold!90!black, rounded corners=4pt, fill=gold!12, align=center, minimum width=2.9cm, minimum height=0.78cm},
    arr/.style={-{Stealth[length=2.4mm]}, thick}
]
\node[box, fill=jade!12, draw=jade] (root) at (0,2) {Trải nghiệm\\du lịch văn hoá};
\node[box] (where) at (-4,0.5) {Nên đi đâu?};
\node[box] (why) at (0,0.5) {Vì sao đáng đi?};
\node[box] (do) at (4,0.5) {Nên làm gì?};
\node[box] (feed) at (-4,-1) {Feed + Map};
\node[box] (facts) at (0,-1) {Tri thức có nguồn};
\node[box] (quest) at (4,-1) {Quest + Capture};
\draw[arr] (root) -- (where);
\draw[arr] (root) -- (why);
\draw[arr] (root) -- (do);
\draw[arr] (where) -- (feed);
\draw[arr] (why) -- (facts);
\draw[arr] (do) -- (quest);
\end{tikzpicture}
\end{frame}

\begin{frame}{7. Input và Output}
\begin{columns}
\column{0.5\textwidth}
\textbf{Input}
\begin{itemize}
    \item Vị trí GPS
    \item Ngôn ngữ, sở thích
    \item Thời tiết, giao thông
    \item Dữ liệu địa điểm
    \item Tri thức văn hoá
    \item Ảnh minh chứng
\end{itemize}
\column{0.5\textwidth}
\textbf{Output}
\begin{itemize}
    \item Feed item
    \item Map marker
    \item Bubble score
    \item Quest task
    \item Capture verdict
    \item Fallback response
\end{itemize}
\end{columns}
\end{frame}

\begin{frame}{8. Pattern Recognition}
\begin{itemize}
    \item Địa điểm nổi tiếng chưa chắc phù hợp khi trời mưa.
    \item Địa điểm gần hơn có thể tốt hơn khi người dùng ít thời gian.
    \item Du khách quốc tế cần giải thích văn hoá, không chỉ bản dịch.
    \item Nhiệm vụ tốt thường có cấu trúc: đi đến nơi, quan sát, chụp ảnh, học thêm.
    \item Thiếu dữ liệu không nên làm app lỗi, mà nên giảm độ tin cậy.
\end{itemize}
\end{frame}

\begin{frame}{9. Abstraction}
\begin{tabularx}{\textwidth}{lX}
\toprule
\textbf{Trừu tượng} & \textbf{Ý nghĩa}\\
\midrule
Place & Địa điểm với danh mục, toạ độ, độ ưu tiên văn hoá.\\
Context & Thời tiết, khoảng cách, giao thông, độ đông, yếu tố thiếu.\\
Cultural Fact & Mẩu tri thức văn hoá có nguồn.\\
Quest Task & Hành động cụ thể có điều kiện hoàn thành.\\
Capture Verdict & Kết quả xác minh ảnh: match, confidence, reason.\\
\bottomrule
\end{tabularx}
\end{frame}

\begin{frame}{10. Thuật toán gợi ý}
\[
S = 0.20W + 0.20T + 0.15D + 0.15C + 0.30K
\]
\begin{itemize}
    \item $W$: điểm thời tiết
    \item $T$: điểm giao thông
    \item $D$: điểm khoảng cách
    \item $C$: điểm độ đông
    \item $K$: điểm văn hoá
\end{itemize}
\begin{block}{Nguyên tắc}
LLM không tự sinh điểm số. Điểm số được tính bằng thuật toán xác định.
\end{block}
\end{frame}

\begin{frame}{11. Điểm văn hoá}
\[
K = 0.40P + 0.30F + 0.20I + 0.10L
\]
\begin{itemize}
    \item $P$: độ ưu tiên nền của địa điểm.
    \item $F$: độ mạnh và độ tin cậy của fact truy xuất.
    \item $I$: mức khớp với sở thích người dùng.
    \item $L$: độ đầy đủ của nội dung và metadata nguồn.
\end{itemize}
\end{frame}

\begin{frame}{12. Dynamic Bubble Map}
\begin{tabular}{lll}
\toprule
\textbf{Điểm phù hợp} & \textbf{Bubble} & \textbf{Ý nghĩa}\\
\midrule
$\geq 70$ & large & Nên ưu tiên hiện tại\\
$45-69$ & medium & Có thể cân nhắc\\
$<45$ & small & Không phải ưu tiên cao\\
\bottomrule
\end{tabular}
\vspace{0.3cm}
\begin{block}{Ý tưởng}
Marker không chỉ là vị trí, mà là biểu diễn trực quan của quyết định tính toán.
\end{block}
\end{frame}

\begin{frame}{13. Storyline Quest}
\begin{enumerate}
    \item Lọc nhiệm vụ theo vị trí hiện tại.
    \item Lọc theo thời tiết và thời điểm.
    \item Bỏ các nhiệm vụ đã hoàn thành.
    \item Chấm điểm nhiệm vụ còn lại.
    \item Trả về nhiệm vụ cụ thể kèm giải thích văn hoá.
\end{enumerate}
\end{frame}

\begin{frame}{14. Capture Judge}
\begin{itemize}
    \item Input: yêu cầu nhiệm vụ + ảnh người dùng upload.
    \item Vision model đánh giá ảnh có khớp nhiệm vụ không.
    \item Output bắt buộc: \texttt{match}, \texttt{confidence}, \texttt{reason}.
    \item Nếu VLM lỗi: fallback trung thực, không tự động approve.
\end{itemize}
\end{frame}

\begin{frame}{15. Kiến trúc Multi-Agent}
\begin{tikzpicture}[
    box/.style={draw=jade, rounded corners=4pt, fill=jade!10, align=center, minimum width=2.7cm, minimum height=0.8cm},
    arr/.style={-{Stealth[length=2.4mm]}, thick}
]
\node[box] (backend) at (-4,0) {Backend};
\node[box] (trip) at (-1,1.2) {Trip Advisor};
\node[box] (culture) at (2,1.2) {Culture Scout};
\node[box] (capture) at (-1,-1.2) {Capture Judge};
\node[box] (vllm) at (2,-1.2) {vLLM};
\draw[arr] (backend) -- (trip);
\draw[arr] (trip) -- (culture);
\draw[arr] (backend) -- (capture);
\draw[arr] (capture) -- (vllm);
\end{tikzpicture}
\vspace{0.25cm}
\begin{block}{Lý do}
Mỗi agent là một phần được phân rã từ bài toán, không phải thêm agent cho phức tạp.
\end{block}
\end{frame}

\begin{frame}{16. Runtime và Offline}
\begin{columns}
\column{0.5\textwidth}
\textbf{Runtime}
\begin{itemize}
    \item Điểm phù hợp hiện tại
    \item Giải thích recommendation
    \item Xác minh capture
\end{itemize}
\column{0.5\textwidth}
\textbf{Offline}
\begin{itemize}
    \item Tri thức đã duyệt
    \item Quest pool
    \item Post giới thiệu địa điểm
    \item Kiểm tra nguồn
\end{itemize}
\end{columns}
\vspace{0.25cm}
Tách như vậy giúp giảm chi phí, độ trễ và hallucination.
\end{frame}

\begin{frame}{17. Độ tin cậy dữ liệu}
\begin{itemize}
    \item Cultural facts phải có source refs.
    \item Nội dung quan trọng nên được chuẩn bị và duyệt trước.
    \item AI chỉ diễn đạt dựa trên dữ liệu đã truy xuất.
    \item Nếu retrieval lỗi, response phải ghi fallback.
\end{itemize}
\end{frame}

\begin{frame}{18. Backend Contract}
\begin{tabularx}{\textwidth}{lX}
\toprule
\textbf{Đối tượng} & \textbf{Field ổn định}\\
\midrule
Feed item & place id, name, score, explanation, source refs\\
Map marker & lat/lng, category, bubble size, score\\
Quest task & title, description, requirement, status\\
Verification & match/status, confidence, reason, fallback\\
\bottomrule
\end{tabularx}
\end{frame}

\begin{frame}{19. Luồng người dùng}
\begin{enumerate}
    \item Người dùng mở app ở Huế.
    \item Feed gợi ý địa điểm gần và phù hợp.
    \item Bản đồ hiển thị bubble theo điểm.
    \item Người dùng xem lý do văn hoá.
    \item Người dùng bắt đầu nhiệm vụ storyline.
    \item Người dùng chụp ảnh minh chứng.
    \item Hệ thống xác minh và mở nội dung tiếp theo.
\end{enumerate}
\end{frame}

\begin{frame}{20. Ví dụ kịch bản}
\begin{block}{Gần Đại Nội Huế, trời nắng}
Địa điểm gần, điểm văn hoá cao, thời tiết phù hợp.
\end{block}
\begin{itemize}
    \item Feed: gợi ý Đại Nội Huế với lý do ``phù hợp để đi hiện tại''.
    \item Map: bubble lớn.
    \item Quest: chụp một chi tiết kiến trúc tại Ngọ Môn.
    \item Giải thích: vai trò nghi lễ và biểu tượng triều Nguyễn.
\end{itemize}
\end{frame}

\begin{frame}{21. Hiện trạng đã làm}
\begin{itemize}
    \item Web app và hướng thiết kế explore map.
    \item Backend endpoints và schema nền tảng.
    \item AI service structure.
    \item Thuật toán recommendation scoring.
    \item Dữ liệu tri thức Huế và hướng quest pool.
    \item Script phục vụ triển khai backend/AI trên VM.
\end{itemize}
\end{frame}

\begin{frame}{22. Phần cần hoàn thiện}
\begin{itemize}
    \item Capture Judge thật bằng vLLM vision endpoint.
    \item LangSmith tracing cho Trip Advisor và Capture Judge.
    \item Seed và kiểm chứng Qdrant retrieval.
    \item Đồng bộ backend contract với web/mobile.
    \item Giảm mock/static data còn lại nếu muốn final product chắc hơn.
\end{itemize}
\end{frame}

\begin{frame}{23. Đánh giá}
\begin{itemize}
    \item Recommendation có giải thích ``vì sao bây giờ'' không?
    \item Bubble size có khớp score không?
    \item Quest có tạo hành động cụ thể không?
    \item Capture verification có verdict đúng schema không?
    \item Khi AI/weather/map lỗi, app có tiếp tục dùng được không?
\end{itemize}
\end{frame}

\begin{frame}{24. Bài học từ tư duy tính toán}
\begin{center}
\Large
Pain point $\rightarrow$ Câu hỏi phân tích $\rightarrow$ Input/Output $\rightarrow$ Abstraction $\rightarrow$ Algorithm $\rightarrow$ Evaluation
\end{center}
\vspace{0.35cm}
\begin{block}{Kết luận}
Đây là cách nhóm biến một vấn đề du lịch văn hoá thành một hệ thống có thể triển khai và đánh giá.
\end{block}
\end{frame}

\begin{frame}{25. Kết thúc}
\begin{center}
\Huge Cảm ơn thầy cô và các bạn
\end{center}
\end{frame}

\end{document}
